\documentclass{article}
\usepackage{graphicx}
\usepackage[T2A]{fontenc}
\usepackage[utf8]{inputenc}
\usepackage[english, russian]{babel}
\usepackage[left=30mm, right=10mm, top=20mm, bottom=20mm]{geometry}
\usepackage{framed}

\usepackage{fancyhdr}
\usepackage{amsmath}

\pagestyle{fancy}
\fancyhf{}
\fancyhead[R]{Шильдт Даниэль гр. 5030102/30202}
\begin{document}
\begin{framed}
    \textbf{I.} Рассмотрим следующую ситуацию: у нас есть некоторое количество монет и мы
    хотим определить, есть ли среди этих монет экземпляры со смещённым центром
    тяжести. Считаем, что у монеты есть смещение, если она выпадает одной и той
    же стороной 10 раз из 10. Примените принцип Бонферрони и укажите при каких
    значениях количества монет это испытание уже не имеет смысла.
\end{framed}
    \textbf{1.} Вероятность того, что честная монета при 10 независимых бросках выпадет одной и той же стороной все 10 раз:\[p_i=P(\text{10 орлов}) + P(\text{10 решек})=\left(\frac{1}{2}\right)^{10}+\left(\frac{1}{2}\right)^{10}=\frac{1}{1024}+\frac{1}{1024}=\frac{1}{512}\approx0.001953125\]\\

    \textbf{2.} Если мы проверим не одну, а $m$ монет, то ожидаемое число "ложных" срабатываний это:\[E[\text{число ложных находок}]=m\cdot p_i\]
    Неформальный принцип Бонферрони гласит: если это ожидаемое число ложных находок существенно больше, чем число реальных аномалий, которые мы надеемся найти, то результат теста бессмысленен, почти все "находки" будут статистическими артефактами, а не настоящими смещёнными монетами.\\

    \textbf{3.} Поскольку в условии не задано конкретное число реально смещенных монет, будем считать, что испытание теряет смысл, когда ожидаемое число ложных срабатываний достигает как минимум одного - то есть тест в среднем находит хотя бы одну ложную аномалию даже среди полностью честных монет\[m\cdot p_i \ge 1\] \[m\cdot \frac{1}{512} \ge 1 \implies m \ge 512\]
    то есть при $m=512$ и более ожидаемое число ложных срабатываний становится $\ge1$, а значит тест теряет практический смысл, мы уже не можем доверять его результатам, потому что даже среди совершенно честных монет в среднем найдется минимум одна, которая случайно "пройдет" тест на смещение.\\
    
    \textbf{Ответ:} при $m \ge 512$ монетах тест теряет практический смысл.

\begin{framed}
    \textbf{II.} Пусть из космоса принимается радиосигнал, который мы интерпретируем как
    последовательность цифр от 0 до 9. Передается по 1 знаку в секунду. Допустим,
    что за время T приема сигнала в этой последовательности были обнаружены 7
    первых знаков числа Пи (3141592). При каких значениях T мы еще можем
    надеяться на разумность отправителя и какова вероятность этого события?
\end{framed}
    \textbf{1.} Вероятность того, что на одной конкретной, заранее выбранной позиции в последовательности окажется ровно 7 нужных цифр в нужной последовательности "3141592": \[P(\text{нужная последовательность})=\left(\frac{1}{10}\right)^{7}=\frac{1}{10\,000\,000}=10^{-7}\]\\

    \textbf{2.} Внутри последовательности из T цифр последовательность из 7 цифр может располагаться $T-6$ способами. Тогда ожидаемое число случайных совпадений во всей последовательности: \[E[\text{число совпадений}]=(T-6)\cdot10^{-7}\]\\

    \textbf{3. } Испытание можно считать значимым, пока ожидаемое число случайных совпадений меньше 1 - то есть найденное совпадение с малой вероятностью объясняется чистой случайностью. Как только $(T-6)\cdot10^{-7}\ge1$, находка теряет доказательную силу. Решая данное неравенство получаем, что при $T\lt10^7+6$ мы еще можем надеяться на разумность отправителя.\\

    \textbf{4.} Позиции внутри последовательности не являются полностью независимыми испытаниями, поэтому точную формулу вероятности "хотя бы одного совпадения" вывести сложно. Однако для оценки можно использовать приближение по аналогии с независимыми испытаниями. В критической точке $T=10^7+6$ имеем $p=10^{-7}$, $n=T-6=10^7$, откуда \[n\cdot p = 10^7\cdot10^{-7}=1.\]Используя приближение $(1-p)^n\approx e^{-np}$ при малом $p$, получаем \[1-(1-p)^n \approx 1-e^{-np}=1-e^{-1}\approx1-0.368=0.632.\]

    \textbf{Ответ:} при $T< 10^7+6$ мы еще можем надеяться на разумность отправителя.\\ В критической точке $T=10^7+6$: вероятность случайного совпадения $\approx0.632$, вероятность того, что находка не случайна (отправитель разумен)$\approx0.368$.
\begin{framed}
    \textbf{III.} Приведите или придумайте примеры ситуаций, где проявляется эффект
    множественных сравнений с применением теста Стьюдента.
\end{framed}
    \textbf{1.} Рассмотрим следующую ситуацию: сравнивается время выполнения одного и того же задания студентами из 10 разных учебных групп. Цель сравнения --- проверить гипотезу о том, влияет ли состав группы, преподаватель или методика обучения на скорость выполнения задания. Если сравнивать каждую группу с каждой другой попарно по критерию Стьюдента, число таких сравнений равно числу сочетаний из 10 по 2:\[\binom{10}{2}=\frac{10\cdot9}{2}=45\]

    \textbf{2.} Пусть вероятность ложного заключения о различии в одном тесте $p_i=0.05$. Тогда вероятность допустить хотя бы одну ошибку среди всех 45 попарных сравнений:\[ p=1-(1-p_i)^{45}=1-(0.95)^{45}\approx0.900\]
    То есть даже если все 10 групп абсолютно одинаковы и реальной разницы во времени выполнения нет, с вероятностью около 90\% при 45 попарных сравнениях найдётся хотя бы одна пара групп с "статистически значимым" различием --- просто по случайности. Отметим, что при $m=45$, что существенно больше 7--8, классическая поправка Бонферрони здесь уже теряет практическую мощность.

    \textbf{3.} Похожий эффект возникает и в других областях. Например, при A/B/C-тестировании нескольких вариантов оформления карточки товара на маркетплейсе: если сравнивать средний чек или конверсию для 6 вариантов оформления попарно ($\binom{6}{2}=15$ сравнений), также существенно вырастает шанс случайно "найти" лучший вариант, который на самом деле не лучше остальных.

\end{document}
